% ---------------------------------------------------------------------------
% INTRODUCTION --- v3 (2026-08-26). Register calibrated against DextrAH-G's
% actual introduction (arXiv:2407.02274v2), fetched verbatim.
%
% Structural moves borrowed from it:
%   P1  vivid hook -> why it matters -> "However, existing methods..." -> "Moreover,"
%   P2  what existing approaches do -> "While effective, ..." -> what is actually needed
%   P3  "Encouragingly, ..." the enabling trend -> "However, ..." what it still misses
%   P4  "To address these challenges, we propose \method: ..." -> contributions as
%       INLINE NUMBERED CLAUSES in a single sentence (not an itemize --- this is
%       what DextrAH-G does and it is much more compact), closing on a
%       forward-looking phrase.
%
% Application-paper ordering: the capability leads, the planner is machinery.
% Claim is deliberately "a framework that supports multi-object gentle grasping",
% NOT "grasp anything" --- it degrades gracefully if the generalist underperforms.
%
% Previous versions in archive/. v1 = pre-literature-sweep; v2 = with citations
% but ML-register and itemized contributions.
% ---------------------------------------------------------------------------
\newcommand{\method}{\textsc{Method}}   % TODO: name the framework

\section{INTRODUCTION}

Handling delicate food is a skill people exercise without thought, yet it remains difficult
for robots. Automating it matters to food processing, agriculture, and assistive settings,
where produce tears, cracks, and bruises under modest compression and the damage is
irreversible~\cite{zhu2025deformable,blanco2024t,wang2025damageless}. A successful grasp must
be right in two respects at once: where and in what orientation the fingers close, and how far
they close. Both depend on the item's pose, size, shape, and stiffness. These properties
vary continuously within a single food category and change entirely across categories, so a
policy that handles one item well need not handle its neighbour. Stiffness is the most awkward
of them to train against. Whatever a robot can perceive of ripeness, the mapping from
appearance to the grip an item needs has to be learned from demonstrations that span the
range, and unlike pose or size, stiffness cannot be dialled in on real hardware: each specimen
contributes one sample at whatever value it happens to have, unmeasured, so covering the range
means sourcing and grading produce rather than simply asking for it.

Existing approaches (see Section~\ref{sec:related}) address this largely through hardware and
sensing. Compliant and pneumatic end-effectors bound contact pressure
mechanically~\cite{wang2024softgripper}, and tactile or force-controlled grippers close the
loop on contact at run time, handling delicate items well~\cite{kang2026learning,vtam}. While
effective, these methods change what the robot can \emph{observe} rather than where its
\emph{supervision} comes from: the policies are learned by imitation from human demonstrations
recorded on real hardware, one item and one pose at a time. Stress-aware grasp planning targets
the damaging quantity directly~\cite{pan2020stressmin}, and finite-element simulation can
evaluate grasp outcomes on deformable objects at scale~\cite{defgraspsim}, but neither has been
used to supply the demonstrations a policy trains on. Acquiring broad coverage of the physical
variation therefore remains a human-labour problem.

Encouragingly, recent work establishes both that coverage is what matters and that it can be
manufactured. Lin \emph{et al.} show that policy generalisation follows a power law in the
number of environments and objects, and that demonstrations per object saturate quickly:
diversity, not count, is what buys generalisation~\cite{lin2025datascaling}. Precision sharpens
the requirement, since for high-precision tasks the number of demonstrations needed grows
super-exponentially as the tolerance tightens, against a limit precision that itself depends on
the quality of the expert~\cite{curseofprecision}. Data-generation systems already multiply a
few hundred human demonstrations into tens of thousands by adapting them to new object
placements~\cite{mandlekar2023mimicgen}. However, such adaptation is geometric, and the
dimension that matters most here cannot be reached by transforming a trajectory: material
stiffness is invisible to a geometric transform, and it cannot be collected on real hardware
either, since one physical specimen cannot be presented at several stiffnesses.

To address these challenges, we propose \method, a framework that generates gentle
demonstrations autonomously in simulation and distils them into a single policy covering many
fragile food categories. Our key observation is that internal stress, the quantity that
determines whether an item is damaged and which no sensor on our robot reports, is computed
exactly inside a deformable simulator; we therefore treat gentleness as a planning objective
rather than a sensed signal. The main contributions of this work include: 1) a stress-aware
grasp planner built on a width-controlled, distributed-pad contact model, whose linearity in
Young's modulus makes randomised material coverage essentially free, 2) an autonomous
demonstration generator that produces gentle, successful grasps across randomised object pose,
size, shape, and material with no human demonstration, 3) a category-conditioned policy trained
on those demonstrations and co-trained with a small slice of real
teleoperation~\cite{simrealcotrain}, acting on a point cloud from a single external depth
camera---a modality we adopt for its narrow appearance gap between simulation and reality
rather than for any perceptual advantage---with no tactile or force sensing in the loop, and
4) real-robot evaluation across
\textbf{[TODO: $N$ categories, headline success rate]} fragile food categories, a step towards
general gentle manipulation of everyday food.
